\documentclass[10pt,a4paper]{article}
\usepackage[T1]{fontenc}
\usepackage[utf8]{inputenc}
\usepackage{lmodern}
\usepackage{amsmath,amssymb,amsthm,mathtools}
\usepackage[margin=0.88in,top=0.92in,bottom=0.88in]{geometry}
\usepackage{microtype}
\usepackage[dvipsnames]{xcolor}
\usepackage{titlesec}
\usepackage{enumitem}
\usepackage{fancyhdr}
\usepackage[hidelinks]{hyperref}

\definecolor{ink}{HTML}{17293E}
\definecolor{accent}{HTML}{275D8C}
\definecolor{muted}{HTML}{536477}
\color{ink}
\titleformat{\section}{\large\bfseries\color{accent}}{\thesection.}{0.55em}{}
\titleformat{\subsection}{\normalsize\bfseries\color{ink}}{\thesubsection}{0.55em}{}
\titlespacing*{\section}{0pt}{1.10em}{0.35em}
\titlespacing*{\subsection}{0pt}{0.75em}{0.22em}
\setlength{\parindent}{0pt}
\setlength{\parskip}{4pt}
\setlist[itemize]{leftmargin=1.35em,topsep=2pt,itemsep=1pt,parsep=0pt}
\setlist[enumerate]{leftmargin=1.5em,topsep=2pt,itemsep=1pt,parsep=0pt}
\setlength{\abovedisplayskip}{7pt}
\setlength{\belowdisplayskip}{7pt}
\setlength{\abovedisplayshortskip}{4pt}
\setlength{\belowdisplayshortskip}{4pt}
\setlength{\headheight}{20pt}
\pagestyle{fancy}
\fancyhf{}
\fancyhead[L]{\footnotesize\color{muted} STAGE 1 \quad / \quad MODULE 1}
\fancyhead[R]{\footnotesize\color{muted} REAL ANALYSIS BASICS}
\fancyfoot[C]{\footnotesize\color{muted}\thepage}
\renewcommand{\headrulewidth}{0.3pt}
\renewcommand{\headrule}{\hbox to\headwidth{\color{muted!35}\leaders\hrule height \headrulewidth\hfill}}
\newcommand{\N}{\mathbb N}
\newcommand{\Z}{\mathbb Z}
\newcommand{\Q}{\mathbb Q}
\newcommand{\R}{\mathbb R}
\newcommand{\Rbar}{\overline{\mathbb R}}
\newcommand{\PO}{\mathcal P(\Omega)}
\newcommand{\ind}[1]{\mathbf 1_{#1}}
\newcommand{\io}{\text{ i.o.}}
\newcommand{\key}[1]{\par\smallskip\textcolor{accent}{\textbf{Key point.}}\ #1\par}

\begin{document}
\begin{center}
  {\LARGE\bfseries Module 1: Foundations and Limit Superior/Inferior}\par
  \vspace{5pt}
  {\small\color{muted}Proof language, sets, countability, sequences, and indicator functions}\par
  \vspace{7pt}
  {\color{accent}\rule{0.88\linewidth}{0.7pt}}
\end{center}
\vspace{-4pt}

\section{Mathematical Statements and Proof Language}
\textbf{Quantifiers.} The symbols $\forall$ (for all) and $\exists$ (there exists) specify the logical scope of a statement. Their order matters:
\[
\forall x\in\R\,\exists y\in\R:\ y>x \quad\text{is true},
\qquad
\exists y\in\R\,\forall x\in\R:\ y>x \quad\text{is false}.
\]
\textbf{Implications.} $P\Rightarrow Q$ means that $P$ is \emph{sufficient} for $Q$, while $Q$ is \emph{necessary} for $P$; $P\Leftrightarrow Q$ means both directions hold.

\textbf{Negation of quantified statements.} Quantifiers switch under negation:
\[
\neg(\forall x\,P(x))\iff\exists x\,\neg P(x),\qquad
\neg(\exists x\,P(x))\iff\forall x\,\neg P(x).
\]
In particular, the distinction between ``eventually'' and ``arbitrarily far along'' is encoded by
\[
\underbrace{\exists N\,\forall n\ge N:\ P(n)}_{\text{eventually always}}
\qquad\text{versus}\qquad
\underbrace{\forall N\,\exists n\ge N:\ P(n)}_{\text{infinitely often}}.
\]
Common proof strategies include direct proof, contrapositive ($P\Rightarrow Q\iff\neg Q\Rightarrow\neg P$), contradiction, and induction.

\section{Sets, Images, and Inverse Images}
For $A,B\subseteq\Omega$, union and intersection correspond to logical OR and AND:
\[
A\cup B=\{x:x\in A\text{ or }x\in B\},\qquad
A\cap B=\{x:x\in A\text{ and }x\in B\}.
\]
The complement is $A^c=\Omega\setminus A$, the difference is $A\setminus B=A\cap B^c$, and the Cartesian product is $A\times B=\{(a,b):a\in A,b\in B\}$. De Morgan's laws extend to any indexed family:
\[
\left(\bigcup_i A_i\right)^c=\bigcap_i A_i^c,
\qquad
\left(\bigcap_i A_i\right)^c=\bigcup_i A_i^c.
\]
For a map $f:X\to Y$, define the \emph{direct image} and \emph{inverse image} by
\[
f[A]=\{f(x):x\in A\}\quad (A\subseteq X),
\qquad
f^{-1}[B]=\{x\in X:f(x)\in B\}\quad (B\subseteq Y).
\]
The preimage $f^{-1}[B]$ exists even if $f$ has no inverse function. Crucially, preimages preserve \emph{all} unions, intersections, and complements:
\[
\boxed{\begin{aligned}
f^{-1}\!\left[\bigcup_i B_i\right]&=\bigcup_i f^{-1}[B_i],\\
f^{-1}\!\left[\bigcap_i B_i\right]&=\bigcap_i f^{-1}[B_i],\\
f^{-1}[Y\setminus B]&=X\setminus f^{-1}[B].
\end{aligned}}
\]
Direct images preserve unions, but generally not intersections or complements. For example, if $f(1)=f(2)=0$, then $f[\{1\}\cap\{2\}]=\varnothing$ whereas $f[\{1\}]\cap f[\{2\}]=\{0\}$.

\key{Inverse images respect the logical structure of subsets. This is why they are fundamental in the definitions of continuity and measurability.}

\section{Countability}
A set is \emph{countably infinite} if it admits a bijection with $\N$; \emph{countable} means finite or countably infinite. Use $\N=\{1,2,3,\ldots\}$ below.

\textbf{$\Z$ is countable.} One enumeration is
\[
0,1,-1,2,-2,3,-3,\ldots,
\]
which lists every integer exactly once, hence defines a bijection $\N\leftrightarrow\Z$.

\textbf{$\Q$ is countable.} Every rational number has a representation $p/q$ with $p\in\Z$ and $q\in\N$. The pairs $(p,q)$ can be enumerated, for example, by increasing $|p|+q$ (only finitely many pairs occur at each level). Enumerate the corresponding fractions and skip repetitions; every rational eventually appears. Since $\Q$ is infinite, this gives a bijection $\N\leftrightarrow\Q$.

\textbf{$\R$ is uncountable (Cantor diagonal argument).} It suffices to show that $(0,1)$ cannot be enumerated. Suppose, toward a contradiction, that
\[
(0,1)=\{x_1,x_2,\ldots\},\qquad x_n=0.a_{n1}a_{n2}a_{n3}\ldots
\]
where each decimal uses the expansion not ending in an infinite string of $9$'s. Construct $y=0.b_1b_2\ldots$ by
\[
b_n=\begin{cases}1,&a_{nn}\ne1,\\2,&a_{nn}=1.\end{cases}
\]
Then $y\in(0,1)$ and its $n$th digit differs from that of $x_n$, so $y\ne x_n$ for every $n$. This contradicts the assumed complete enumeration.

\key{For $\Z$ and $\Q$, \emph{one} bijection proves countability. For $(0,1)$, diagonalization defeats \emph{every} proposed enumeration.}

\section{Sequences and Their Limit Superior/Inferior}
\subsection{Convergence and bounds (brief review)}
A real sequence converges to $L$ if $\forall\varepsilon>0\,\exists N\,\forall n\ge N$, $|x_n-L|<\varepsilon$. For a nonempty bounded set $S\subseteq\R$, $\sup S$ is its least upper bound and $\inf S$ its greatest lower bound; neither need belong to $S$. The least-upper-bound property of $\R$ is a formulation of completeness.

\subsection{Numerical sequences}
For $(x_n)\subseteq\R$, define the upper and lower bounds of the $N$th tail by
\[
s_N=\sup_{n\ge N}x_n,\qquad t_N=\inf_{n\ge N}x_n.
\]
Deleting initial terms can only decrease $s_N$ and increase $t_N$:
\[
s_{N+1}\le s_N,\qquad t_{N+1}\ge t_N.
\]
Thus both auxiliary sequences have limits in the \emph{extended reals} $\Rbar=\R\cup\{-\infty,+\infty\}$, and we define
\[
\boxed{\begin{aligned}
\limsup_{n\to\infty}x_n&=\lim_{N\to\infty}s_N=\inf_N\sup_{n\ge N}x_n,\\
\liminf_{n\to\infty}x_n&=\lim_{N\to\infty}t_N=\sup_N\inf_{n\ge N}x_n.
\end{aligned}}
\]
The outer $\inf$ and $\sup$ arise because the tail upper bounds decrease and the tail lower bounds increase. They are \emph{asymptotic} bounds, not necessarily the bounds of the whole sequence.

For a bounded sequence, $\limsup x_n$ and $\liminf x_n$ are respectively the largest and smallest \emph{subsequential limits}. The limit values need not be attained by any term. In addition,
\[
\liminf x_n\le\limsup x_n,\qquad
x_n\to L\in\R\iff\liminf x_n=\limsup x_n=L.
\]
For instance, $x_n=(-1)^n+1/n$ satisfies $\limsup x_n=1$ and $\liminf x_n=-1$; an isolated first term of $100$ would not affect either limit. An unbounded sequence may have an infinite upper or lower limit (e.g., $x_n=(-1)^n n$).

\subsection{Sequences of sets}
Let $A_n\subseteq\Omega$. In the inclusion order on $\PO$, the supremum of a family is its union and the infimum is its intersection. Define the tail upper and lower bounds
\[
B_N=\sup_{n\ge N}A_n=\bigcup_{n\ge N}A_n,\qquad
C_N=\inf_{n\ge N}A_n=\bigcap_{n\ge N}A_n.
\]
As $N$ increases, $B_N$ decreases under inclusion while $C_N$ increases. Consequently,
\[
\boxed{\begin{aligned}
\limsup_{n\to\infty}A_n&=\inf_N B_N=\bigcap_{N=1}^{\infty}\bigcup_{n\ge N}A_n,\\
\liminf_{n\to\infty}A_n&=\sup_N C_N=\bigcup_{N=1}^{\infty}\bigcap_{n\ge N}A_n.
\end{aligned}}
\]
Equivalently, expanding the membership quantifiers,
\[
\begin{aligned}
x\in\limsup A_n&\iff \forall N\,\exists n\ge N:\ x\in A_n
&&\text{(infinitely often)},\\
x\in\liminf A_n&\iff \exists N\,\forall n\ge N:\ x\in A_n
&&\text{(eventually always)}.
\end{aligned}
\]
These descriptions are \emph{consequences} of the order-theoretic definitions: the outer operation combines all tails, not just one fixed tail.

\section{The Order-Theoretic Bridge: Indicator Functions}
The power set $(\PO,\subseteq)$ is a \emph{complete lattice}: every family of subsets has a supremum (union) and an infimum (intersection). It is therefore meaningful to use the same tail sup/inf construction for numbers and sets.

For $A\subseteq\Omega$, its \emph{indicator function} is
\[
\ind A(x)=\begin{cases}1,&x\in A,\\0,&x\notin A.\end{cases}
\qquad\text{Thus}\qquad
A=\{x:\ind A(x)=1\}.
\]
It translates set membership into a numerical value, without losing information. In fact $A\mapsto\ind A$ is an \emph{order isomorphism} from $(\PO,\subseteq)$ to the $\{0,1\}$-valued functions on $\Omega$ with pointwise order:
\[
\begin{gathered}
A\subseteq B\iff\ind A\le\ind B,\qquad
\ind{A^c}=1-\ind A,\qquad
\ind{A\cap B}=\ind A\,\ind B,\\
\ind{\bigcup_i A_i}(x)=\sup_i\ind{A_i}(x),\qquad
\ind{\bigcap_i A_i}(x)=\inf_i\ind{A_i}(x).
\end{gathered}
\]
Apply these identities successively to the two layers of the set limit:
\[
\begin{aligned}
\ind{\limsup A_n}(x)
&=\ind{\bigcap_N\bigcup_{n\ge N}A_n}(x)
=\inf_N\sup_{n\ge N}\ind{A_n}(x)
=\limsup_n\ind{A_n}(x),\\
\ind{\liminf A_n}(x)
&=\sup_N\inf_{n\ge N}\ind{A_n}(x)
=\liminf_n\ind{A_n}(x).
\end{aligned}
\]
For a fixed $x$, the sequence $\ind{A_n}(x)$ takes values only in $\{0,1\}$. Its limsup equals $1$ exactly when $x$ appears infinitely often, while its liminf equals $1$ exactly when $x$ eventually appears at every index.

Moreover, $\liminf A_n\subseteq\limsup A_n$: for any $m,N$, choosing $k=\max\{m,N\}$ gives
\[
C_m\subseteq A_k\subseteq B_N
\quad\Longrightarrow\quad
\sup_m C_m\subseteq\inf_N B_N.
\]
Finally, convergence of sets in this sense is exactly pointwise convergence of indicators:
\[
\boxed{\liminf A_n=\limsup A_n=A
\iff\ind{A_n}(x)\to\ind A(x)\quad\forall x\in\Omega.}
\]
\key{The indicator map turns a problem about sets into one about numerical functions. Later, measure and probability will use the same bridge: $\mu(A)=\int\ind A\,d\mu$ and $\mathbb P(A)=\mathbb E[\ind A]$.}

\end{document}
